% META: {"id": "TSPE-PRIMEQ-ALG-015", "chapitre": "TSPE-PRIMITIVES-EQDIFF", "type_objet": "algorithme", "capacites_codes": ["C1", "C2"], "programme_alignment": "OFFICIAL_ALGORITHM_EXAMPLE", "mode_creation": "ex_nihilo", "sources_inspiration": [], "status": "generated"}

\section{Algorithmique : résolution approchée par la méthode d'Euler}

% ============================================================
% STRATE 1 --- Probleme, modele, algorithme
% ============================================================

\margeAppui{
« Resolution par la methode d'Euler de $y' = f$, de $y' = ay + b$ » figure
parmi les \textbf{exemples d'algorithme} du programme. Il la cite, il ne
l'impose pas.
}

\textbf{Le probleme.} On sait resoudre $y' = ay + b$ exactement. C'est
precisement pour cela que ce cas est le bon terrain d'essai : il permet de
\emph{mesurer} l'erreur d'une methode approchee, avant de l'employer la ou
aucune solution explicite n'existe.

\propriete[Le schema d'Euler]{
Sur un petit pas $h$, on remplace la courbe par sa tangente :
$y(t+h) \approx y(t) + h\,y'(t)$. Pour l'equation $y' = ay + b$ avec la
condition initiale $y(0) = y_0$, cela donne
\[
  t_{k+1} = t_k + h,
  \qquad
  y_{k+1} = y_k + h\left(a\,y_k + b\right).
\]
Entree : $a$, $b$, $y_0$, le pas $h$ et la duree. Sortie : la liste des points
$(t_k\,;\,y_k)$.
}

\begin{python}
def euler_lineaire(a, b, y0, pas, duree):
    """Solution approchee de y' = a y + b sur [0 ; duree], par Euler."""
    if pas <= 0 or duree < 0:
        raise ValueError("le pas doit etre strictement positif")
    t, y = 0.0, y0
    points = [(t, y)]
    while t < duree - pas / 2:
        y = y + pas * (a * y + b)
        t = t + pas
        points.append((t, y))
    return points
\end{python}

\exemple{
Prenons $y' = -2y + 6$ avec $y(0) = 1$. La solution exacte est
\[
  y(t) = 3 - 2\,\mathrm{e}^{-2t},
  \qquad\text{d'ou}\qquad
  y(1) = 3 - \frac{2}{\mathrm{e}^{2}} = 2{,}729\,329\ldots
\]

\begin{center}
\begin{tabular}{|c|c|c|}
\hline
pas $h$ & Euler en $t=1$ & ecart a la solution \\
\hline
$0{,}1$ & $2{,}785\,252$ & $-0{,}055\,922$ \\
\hline
$0{,}05$ & $2{,}756\,847$ & $-0{,}027\,517$ \\
\hline
$0{,}01$ & $2{,}734\,761$ & $-0{,}005\,431$ \\
\hline
$0{,}001$ & $2{,}729\,871$ & $-0{,}000\,542$ \\
\hline
\end{tabular}
\end{center}

Deux constats. L'ecart est ici toujours \textbf{negatif} : Euler surestime,
parce que la solution est concave et que la tangente passe au-dessus. Et
diviser le pas par dix divise l'ecart par dix : l'erreur de la methode d'Euler
est proportionnelle au pas -- on dit qu'elle est d'ordre $1$.
}

\propriete[Travail demande]{%
\begin{enumerate}
  \item Retrouver a la main $y_1$ et $y_2$ pour $h = 0{,}1$.
  \item Verifier que la solution exacte satisfait bien $y(0) = 1$ et
        $y' = -2y+6$.
  \item Estimer le pas qui donnerait quatre decimales exactes en $t = 1$.
  \item Reprendre avec $a = 2$, $b = 0$, $y_0 = 1$ : l'ecart change-t-il de
        signe ? Expliquer avec la convexite.
\end{enumerate}
}

\exemple{
\textbf{Reponses.} 3. L'erreur etant proportionnelle au pas et valant
$5{,}4\times10^{-3}$ pour $h = 10^{-2}$, il faut $h$ de l'ordre de
$10^{-4}$. 4. Oui : la solution $y = \mathrm{e}^{2t}$ est \emph{convexe}, sa
courbe est au-dessus de ses tangentes, et Euler sous-estime.
}

% BEGIN-VERIFY
% from math import exp
%
% def euler_lineaire(a, b, y0, pas, duree):
%     if pas <= 0 or duree < 0:
%         raise ValueError("le pas doit etre strictement positif")
%     t, y = 0.0, y0
%     points = [(t, y)]
%     while t < duree - pas / 2:
%         y = y + pas * (a * y + b)
%         t = t + pas
%         points.append((t, y))
%     return points
%
% exacte = lambda t: 3 - 2 * exp(-2 * t)
% assert exacte(0) == 1
% assert abs(exacte(1) - 2.729329) < 5e-7
% # La solution annoncee verifie bien l'equation : y' = -2y + 6.
% for t in (0.0, 0.3, 1.0, 2.5):
%     derivee_numerique = (exacte(t + 1e-7) - exacte(t - 1e-7)) / 2e-7
%     assert abs(derivee_numerique - (-2 * exacte(t) + 6)) < 1e-5, t
%
% valeurs = {}
% for h, attendue, ecart_attendu in (
%     (0.1, 2.785252, -0.055922),
%     (0.05, 2.756847, -0.027517),
%     (0.01, 2.734761, -0.005431),
%     (0.001, 2.729871, -0.000542),
% ):
%     y = euler_lineaire(-2.0, 6.0, 1.0, h, 1.0)[-1][1]
%     valeurs[h] = y
%     assert round(y, 6) == attendue, (h, y)
%     assert round(exacte(1) - y, 6) == ecart_attendu, (h, exacte(1) - y)
%     # Euler surestime : la solution est concave.
%     assert y > exacte(1)
%
% # Erreur d'ordre 1 : diviser le pas par dix divise l'ecart par dix.
% assert 9 < (exacte(1) - valeurs[0.01]) / (exacte(1) - valeurs[0.001]) < 11
% # Quatre decimales demandent un pas de l'ordre de 10^-4.
% fin = euler_lineaire(-2.0, 6.0, 1.0, 1e-4, 1.0)[-1][1]
% assert abs(exacte(1) - fin) < 1e-4
%
% # Les deux premiers pas, a la main : y1 = 1 + 0,1*(-2+6) = 1,4.
% debut = euler_lineaire(-2.0, 6.0, 1.0, 0.1, 0.2)
% assert abs(debut[1][1] - 1.4) < 1e-12
% assert abs(debut[2][1] - (1.4 + 0.1 * (-2 * 1.4 + 6))) < 1e-12
%
% # Solution convexe : l'ecart change de signe.
% croissante = euler_lineaire(2.0, 0.0, 1.0, 0.01, 1.0)[-1][1]
% assert croissante < exp(2)
%
% for mauvais in ((-2.0, 6.0, 1.0, 0.0, 1.0), (-2.0, 6.0, 1.0, -0.1, 1.0)):
%     try:
%         euler_lineaire(*mauvais)
%     except ValueError:
%         pass
%     else:
%         raise AssertionError("le pas invalide aurait du etre refuse")
% END-VERIFY

% ============================================================
% STRATE 2 --- Erreurs frequentes
% ============================================================

\erreurFrequente{
\textbf{Croire qu'Euler donne la solution.} Il en donne une valeur approchee,
avec une erreur proportionnelle au pas. Sur cet exemple, dix pas de $0{,}1$
se trompent des la deuxieme decimale.
}

\erreurFrequente{
\textbf{Oublier que le sens de l'erreur depend de la convexite.} La tangente
est au-dessus d'une courbe concave et au-dessous d'une courbe convexe : Euler
surestime dans le premier cas, sous-estime dans le second. Le signe de l'ecart
n'est pas une propriete de la methode, mais de la solution.
}
